\section{Discussion}
\subsection{An Executable Interface Is Distinct from a Better Tradeoff}
The cumulative budget provides a common unit linking terminal supply forecasts, stop-level legal commands and physical execution. Its accounting identity answers whether the same limited resource was used throughout a trajectory. The independent retention study answers another question: how the operating point moved when that resource constraint changed. Both have positive evidence here. Neither by itself establishes that pressure-driven allocation outperforms a simpler restriction at the same service quality.

The independent simple-control comparison supplies that missing challenge. Pressure credit reduced dispatch delay but failed the registered regularity protections, while its journey benefit depended on the comparator. This is consistent with a useful operating configuration rather than a uniformly improved tradeoff. It also explains why holding exposure must be reported: matching average exposure on old days does not fix the exposure, service states or regularity of new trajectories.

This distinction is consistent with the established need to evaluate headway stability alongside holding time~\cite{berrebi2018comparing}. Our addition is the explicit accounting of published terminal obligations and a persistent vehicle--trip allowance under a fixed pool. The results identify the present boundary of that interface rather than a general resolution of the holding tradeoff.

\subsection{The Protected Reference Determines the Question}
The original fast selector protects a reference that has already been changed by the budget. Thus even an accurate local comparison can preserve the budgeted continuation while the budget itself removes useful compensation. The common-reference redesign changed the question to whether a new restriction and its current action improved upon the no-new-tightening continuation.

That redesign had the intended directional effect on actual regularity, but it admitted few restrictions and lost the operating savings. Its recorded behavior suggests that the current common-reference conditions were restrictive in these states. The data do not isolate whether a particular CV$^2$, resource-timing or passenger-cost condition was primarily responsible for each rejected request. We therefore interpret it as another measured operating point, not as proof that forecasts alone caused the failure or that a looser retrospective threshold would succeed.

\subsection{Learning Should Be Judged After Matching the Interface}
The same-budget and explicit-retention studies constrain the interpretation of the learned component. At the default setting there was no established journey increment over manual joint; at the tighter setting the independent estimate favored manual joint. Explicit budget context did not change the final executed mode. The presence of training, active timetable revisions or differing individual commands is insufficient to infer an improvement in daily passenger outcomes.

The harmonic comparison also separates representation from full frequency specialization. Retaining harmonics favored the secondary regularity endpoint but failed the primary journey criterion on new days. The fast controller continued to use APC clearance forecasts and AVL rather than a direct harmonic-residual command. Hence the current method is a frequency-informed coordination architecture with measured representation effects, not an empirical demonstration of a complete low-/high-frequency division of labor.

\subsection{Limitations and Use}
\label{sec:limitations}
The independent panels cover new realizations of one frozen simulation-demand generator on one corridor, fleet size and trip inventory. They are not route-family transfer, field deployment, or an estimate of causal effects on observed passengers. The 16 test days concern one frozen learner seed; they do not establish robustness across independent learning initializations. The four-day studies remain development screens, and the archived gates are not replaced by favorable secondary endpoints.

The graph forecasts demand, service timing and native continuation under the current published state. Future fast replanning, later credit publications and uncertain running conditions can change the realized trajectory. Its relative constraints and cumulative-credit identity therefore do not establish a complete-day performance guarantee. Holding has only nonnegative legal actions, and the bounded timetable law acts within a fixed fleet; the studies do not identify the best result under a different vehicle inventory or dispatch authority.

The simple controls are prespecified daily mixtures calibrated on mean development exposure. Their expected daily endpoints are valid for those randomized settings, but do not represent within-trajectory interpolation or a test-day matched-dose experiment. The comparison concerns these implemented simple families, not all possible static restrictions. Likewise, readiness and public-delay improvements share operational causes and should not be interpreted as independent evidence sources.

The budget parameter has an explicit command-time meaning, but expert usability, dispatcher workload and acceptable service-quality tolerances were not measured. No tolerance is inferred retrospectively from the observed CV$^2$ increase. Full native runtime portability and a public data-release arrangement are separate from regenerating the paper tables. The current evidence is most useful as a reproducible account of configuration and comparison choices; a stronger claim requires an independently tested improvement over the simple-control tradeoff.
